\documentclass[openany]{book}

\usepackage[margin=1in]{geometry}
\usepackage{amsmath,amsfonts,amsthm, amssymb}
\usepackage{yhmath}
\usepackage{mathrsfs}
\usepackage{mathtools}
\usepackage{xcolor}
\usepackage{graphicx}
\usepackage{comment}
\usepackage{tikz-cd}
\usepackage{quiver}
\usepackage{hyperref,cleveref}
\renewcommand{\familydefault}{ppl}
\newcommand{\curl}{\text{curl}}
\newcommand{\diverg}{\text{div}}
\newcommand{\tr}{\text{tr}}
\newcommand{\R}{\mathbb{R}}
\newcommand{\E}{\mathbb{E}}
\newcommand{\Z}{\mathbb{Z}}
\newcommand{\C}{\mathbb{C}}
\newcommand{\F}{\mathbb{F}}
\newcommand{\la}{\langle}
\newcommand{\ra}{\rangle}
\newcommand{\colim}{\text{colim}}
\DeclareMathOperator{\im}{im}
\DeclareMathOperator{\disc}{disc}
\let\oldemptyset\emptyset
\let\emptyset\varnothing
\newcommand{\tor}{\text{Tor}}
\newcommand{\id}{\text{id}}
\newcommand{\ext}{\text{Ext}}
\newcommand{\ptop}{\text{PTop}}
\newcommand{\pt}{\text{pt}}
\newcommand{\ach}{\text{Ach}}
\newcommand{\Q}{\mathbb{Q}}
\newcommand{\gal}{\text{Gal}}
\newcommand{\fraccomma}{\genfrac{}{}{0pt}{}{}{,}}
\definecolor{wikipediadarkblue}{rgb}{0.023, 0.270, 0.676}
\hypersetup{
    colorlinks,
    citecolor=black,
    filecolor=black,
    linkcolor=blue,
    urlcolor=red
}

% \hypersetup{
%     urlbordercolor={1 0 0}, % Red box for URLs
%     linkbordercolor=blue, % Blue box for internal links
%     % the values are in RGB format from 0 to 1
% }


\input{hui_r.tex}

\title{Calc III Midterm 2 Review
\\ 
\vspace{0.4cm}
\large Spring 2026}



\author{(Please email hsun95@jh.edu if you see typos)}
\date{\today}


\begin{document}

\maketitle

\tableofcontents
\newpage


% \chapter{Definition Review}


% \begin{defn}[standard basis in $\R^3$]
%     The vectors 
%     \begin{equation*}
%         i=(1,0,0), j=(0,1,0), k=(0,0,1)
%     \end{equation*}
%     are called the \textbf{standard basis} vectors of $\R^3$, and for any vector $a=(a_1,a_2,a_3)\in\R^3$, we can write 
%     \begin{equation*}
%         a=a_1i+a_2j+a_3k
%     \end{equation*}
% \end{defn}




% \begin{defn}[Equation of a line]\label{line}
%     A \textbf{line} $l$ in $\R^3$ through the tip of $a=(a_1,a_2,a_3)$ pointing in the direction of a vector $v=(v_1,v_2,v_3)$ is given by 
%     \begin{equation*}
%         l(t)=a+tv=(a_1+tv_1, a_2+tv_2, a_3+tv_3)
%     \end{equation*}
%     where $t\in\R$. Alternatively, a line passing through two points $P=(x_1,y_1,z_1), Q=(x_2,y_2,z_2)$ is given by 
%     \begin{equation*}
%         l(t)=(x(t), y(t), z(t))
%     \end{equation*}
%     where 
%     \begin{equation*}
%         \begin{cases}
%             x(t)=x_1+(x_2-x_1)t\\
%             y(t)=y_1+(y_2-y_1)t\\
%             z(t)=z_1+(z_2-z_1)t
%         \end{cases}
%     \end{equation*}
% \end{defn}



% \begin{defn}[inner product, dot product]
%     Let $a,b\in\R^3$, the \textbf{dot product}, also called the inner product, of $a,b$ is 
%     \begin{equation*}
%         a\cdot b=a_1b_1+a_2b_2+a_3b_3
%     \end{equation*}
%     where $a=(a_1,a_2,a_3), b=(b_1,b_2,b_3)$. The \textbf{norm}, also called the length, of $a$ is 
%     \begin{equation*}
%         \|a\|=(a\cdot a)^\frac{1}{2}
%     \end{equation*}
%     A vector of norm $1$ is called a \textbf{unit vector}. Given any $u\in\R^3$, we can find a unit vector $\frac{u}{\|u\|}$ pointing in the same direction as $u$, this is called ``normalizing'' $u$.
% \end{defn}


% \begin{defn}[orthogonal projection]
%     The \textbf{orthogonal projection} of vector $v$ onto another vector $a$ is
%     \begin{equation*}
%         \text{Proj}_av=\frac{a\cdot v}{a\cdot a}a
%     \end{equation*}
%     For example, the orthogonal projection of $(1,1,0)$ onto $(1,1,1)$ is 
%     \begin{equation*}
%         \left(\frac{2}{3}, \frac{2}{3},\frac{2}{3}\right)
%     \end{equation*}
% \end{defn}


% \begin{defn}[orthogonal]
%     Let $a,b\in\R^n$, then $a,b$ are called \textbf{orthogonal} or perpendicular iff 
%     \begin{equation*}
%         a\cdot b=0
%     \end{equation*}
% \end{defn}


% \begin{defn}[determinant]
%     The \textbf{determinant} of a $2\times 2$ matrix is given by 
%     \begin{equation*}
%         \det\begin{bmatrix}
%             a&b\\
%             c&d
%         \end{bmatrix}=ad-bc
%     \end{equation*}
%     and the determinant of a $3\times 3$ matrix is given by 
%     \begin{equation*}
%         \det\begin{bmatrix}
%             a_{11}&a_{12}&a_{13}\\
%             a_{21}&a_{22}&a_{23}\\
%             a_{31}&a_{32}&a_{33}\\
%         \end{bmatrix}=a_{11}(a_{22}a_{33}-a_{23}a_{32})-a_{12}(a_{21}a_{33}-a_{23}a_{31})+a_{13}(a_{21}a_{32}-a_{22}a_{31})
%     \end{equation*}
% \end{defn}



% \begin{defn}[cross product]
%     Let $a,b\in\R^3$, write $a=(a_1,a_2,a_3), b=(b_1,b_2,b_3)$, then the \textbf{cross product}
%     \begin{equation*}
%         a\times b=\det\begin{bmatrix}
%             i&j&k\\
%             a_1&a_2&a_3\\
%             b_1&b_2&b_3
%         \end{bmatrix}
%     \end{equation*}
%     where $i,j,k$ are the standard vectors in $\R^3$.
% \end{defn}

% \begin{defn}[Plane in $\R^3$]\label{plane}
%    If a plane ${P}$ passes through some point $(x_0,y_0,z_0)$, and $n=(A,B,C)$ is a vector orthogonal to the plane, then the plane ${P}$ is given by the equation:
%     \begin{equation*}
%         A(x-x_0)+B(y-y_0)+C(z-z_0)=0
%     \end{equation*}
%     (Notice that a point in ${P}$ and a normal vector to ${P}$ uniquely define a plane in $\R^3$.)
% \end{defn}


% \begin{defn}[image, graph]
%     The \textbf{image} of a function $f: U\subset\R^n\to\R^m$ is a subset of $\R^m$,
%     \begin{equation*}
%         \text{Image}(f)=\{f(x)\in\R^m: x\in U\}
%     \end{equation*}
%     and the \textbf{graph} of $f$ is a subset of $\R^{n+m}$,
%     \begin{equation*}
%         \text{Graph}(f)=\{(x,f(x)): x\in U\}
%     \end{equation*}
% \end{defn}


% \begin{defn}[level set]
%     Let $f:U\subset\R^n\to\R^m$, and $c\in\R$ be some constant. Then the \textbf{level set} of $f$ at $c$ is the set 
%     \begin{equation*}
%         \{x\in U: f(x)=c\}\subset\R^n
%     \end{equation*}
% \end{defn}

% \begin{defn}[open set, closed set, neighborhood, boundary]
%     Let $U\subset\R^n$, we say $U$ is an \textbf{open set} if for every $x_0\in U$, there exists some $r>0$ such that $D_r(x_0)\subset U$, where $D_r(x_0)$ is the open ball of radius $r$ centered at $x_0$:
%     \begin{equation*}
%         D_r(x_0)=\{x\in\R^n: \|x-x_0\|<r\}
%     \end{equation*}
%     Some examples of open sets: $\R$, $D_{1}((0,0))$, $(1,2)\subset\R$.
%     A \textbf{neighborhood} of $x_0\in\R^n$ is an open set containing the point $x_0$. A point $x\in\R^n$ is called a \textbf{boundary point} of $A$ if \textit{every} neighborhood of $x$ contains at least one point in $A$ and at least one point not in $A$. A set is \textbf{closed} if it contains all its boundary points. Example of closed set: level sets of a continuous function $f$.
% \end{defn}


% \begin{defn}[limit]
%     Let $f: A\subset\R^n\to\R^m$, where $A$ is open, let $x_0$ be in $A$ or be a boundary point of $A$ and $N$ be a neighborhood of a point $b\in\R^m$. Now let $x$ approach $x_0$, $f$ is said to be \textbf{eventually in $N$} if there exists a neighborhood $U$ of $x_0$ such that 
%     \begin{equation*}
%         \text{ if } x\in U, \text{ then } f(x)\in N
%     \end{equation*}
%     If $f$ is eventually in $N$ for \textit{any} neighborhood $N$ around $b$, then the \textbf{limit} of $f$ as $x\to x_0$ exists, denoted as 
%     \begin{equation*}
%         \lim_{x\to x_0}f(x)=b
%     \end{equation*}
% \end{defn}

% \begin{defn}[continuous]
%     Let $f:A\subset\R^n\to\R^m$ and $x_0\in A$, then $f$ is \textbf{continuous at} \boldmath{${x_0}$} \unboldmath if 
%     \begin{equation*}
%         \lim_{x\to x_0}f(x)=f(x_0)
%     \end{equation*}
% \end{defn}


% \begin{defn}[partial derivative]
%     Let $f:U\subset\R^n\to\R$, where $U$ is open. Then the \textbf{partial derivative} with respect to $x_i$ is defined by 
%     \begin{equation*}
%         \frac{\partial f}{\partial x_i}(x_1, \dots, x_n)=\lim_{h\to 0}\frac{f(x+he_i)-f(x)}{h}
%     \end{equation*}
%     where $e_i=(0,\dots, 1,\dots, 0)$ with $1$ in the $i$th coordinate. 
% \end{defn}



% \begin{defn}[differentiability in two variables]\label{tangent plane of graph}
%     Let $f:\R^2\to\R$, then $f$ is \textbf{differentiable} at $(x_0,y_0)$ if 
%     \begin{enumerate}
%         \item[(1)] $\frac{\partial f}{\partial x},\frac{\partial f}{\partial y}$ exist at $(x_0,y_0)$
%         \item[(2)] 
%         \begin{equation*}
%             \lim_{(x,y)\to(x_0,y_0)}\frac{f(x,y)-f(x_0,y_0)-\left[\frac{\partial f}{\partial x}(x_0,y_0)\right](x-x_0)-\left[\frac{\partial f}{\partial y}(x_0,y_0)\right](y-y_0)}{\|(x,y)-(x_0,y_0)\|}=0
%         \end{equation*}
%     \end{enumerate}
%     The derivative of $f$ at $(x_0,y_0)$ is the $1\times 2$ matrix 
%     \begin{equation*}
%         \begin{bmatrix}\frac{\partial f}{\partial x}(x_0,y_0)&\frac{\partial f}{\partial y}(x_0,y_0)\end{bmatrix}
%     \end{equation*}
%     Moreover, the \textbf{tangent plane} of the graph of $f$ at $(x_0,y_0, f(x_0,y_0))$ is given by 
%     \begin{equation*}
%         z=f(x_0,y_0)+\left[\frac{\partial f}{\partial x}(x_0,y_0)\right](x-x_0)+\left[\frac{\partial f}{\partial y}(x_0,y_0)\right](y-y_0)
%     \end{equation*}
% \end{defn}

% \begin{defn}[differentiability in the general setting]
%     Let $f:U\subset\R^n\to\R^m$, then $f$ is differentiable at $x_0\in U$ if
%     \begin{enumerate}
%         \item[(1)] the partial derivatives $\frac{\partial f_i}{\partial x_j}$ exist for all $1\leq i\leq m, 1\leq j\leq n$. 
%         \item[(2)]
%         \begin{equation*}
%             \lim_{x\to x_0}\frac{\|f(x)-f(x_0)-T(x-x_0)\|}{\|x-x_0\|}=0
%         \end{equation*}
%         where $T=Df(x_0)$ is the $m\times n$ matrix 
%         \begin{equation*}
%             Df(x_0)=
%             \begin{bmatrix}
%             \frac{\partial f_1}{\partial x_1}(x_0) & \cdots & \frac{\partial f_1}{\partial x_n}(x_0) \\
%             \vdots & \ddots & \vdots \\
%             \frac{\partial f_m}{\partial x_1}(x_0) & \cdots & \frac{\partial f_m}{\partial x_n}(x_0)
%             \end{bmatrix}
%         \end{equation*}
%     \end{enumerate}
%     The derivative of $f$ at $x_0$ is the $m\times n$ matrix $Df(x_0)$.
% \end{defn}

% \begin{defn}[gradient]
%     Let $f:U\subset\R^n\to\R$, the \textbf{gradient} $\nabla f(x)$ is a special case of the general case above when $m=1$, i.e., it is a $1\times n$ matrix 
%     \begin{equation*}
%         Df(x)=\begin{bmatrix}
%             \frac{\partial f}{\partial x_1}&\dots&\frac{\partial f}{\partial x_n}
%         \end{bmatrix}
%     \end{equation*}
% \end{defn}


% \begin{defn}[path and curve]
%     A \textbf{path} in $\R^n$ is a map $c:[a,b]\to\R^n$, and the image of $c$ is called a \textbf{curve}. We say the path $c$ parametrizes the curve. 

%     \noindent For example, $c(t)=(\cos t, \sin t)$ is a path, and the unit circle is a curve.
% \end{defn}


% \begin{defn}[velocity of a path]
%     Let $c:[a,b]\to\R^n$ be a path, and we can write $c(t)=(c_1(t), \dots, c_n(t))$. If $c$ is differentiable, then we define the \textbf{velocity} of $c$ at any $t_0\in [a,b]$ as 
%     \begin{equation*}
%         c'(t_0)=\left(c_1'(t_0), \dots, c_n'(t_0)\right)
%     \end{equation*}
%     The velocity vector of $c$ at $t_0$ is also a \textbf{tangent} vector to $c$ at $t_0$. The \textbf{speed} of the path $c$ at $t_0$ is the length of the velocity vector $\|c'(t_0)\|$.
% \end{defn}

% \begin{defn}[tangent line to a path]\label{tangent path}
%     Let $c:[a,b]\to\R^n$ be a path, if $c'(t_0)\neq 0$, then the \textbf{tangent line} at $x_0$ is given by 
%     \begin{equation*}
%         l(t)=c(t_0)+c'(t_0)(t-t_0)
%     \end{equation*}
% \end{defn}




% \begin{defn}[directional derivative]
%     Let $f:\R^3\to\R$, be differentiable, then the \textbf{directional derivative} at $x_0\in\R^3$ in the direction of a \textit{unit vector} $v$ is given by 
%     \begin{equation*}
%         \nabla f(x_0)\cdot v=\left[\frac{\partial f}{\partial x_1}(x_0)\right]v_1+\left[\frac{\partial f}{\partial x_2}(x_0)\right]v_2+\left[\frac{\partial f}{\partial x_3}(x_0)\right]v_3
%     \end{equation*}
%     where $v=(v_1,v_2,v_3)$.
% \end{defn}
% \begin{warn}
%     Make sure you normalize any given direction $v$! This formula works for unit vectors.
% \end{warn}


\chapter{Definition Review}

\begin{defn}[First order Taylor expansion]
    Let $f:U\subset\R^n\to\R$ be differentiable at $a\in U$, then 
    \begin{equation*}
        f(x)=f(a)+\sum_{i=1}^n\frac{\partial f}{\partial x_i}(a)(x_i-a_i)+R_1(a,x)
    \end{equation*}
    where 
    \begin{equation*}
        \frac{R_1(a,x)}{\|x-a\|}\to 0 \text{ as } x\to a
    \end{equation*}
\end{defn}

\begin{defn}[Second order Taylor expansion]\label{taylor}
    Let $f:U\subset\R^n\to\R$ be twice continuously differentiable at $a\in U$, then 
    \begin{equation*}
        f(x)=f(a)+\sum_{i=1}^n\frac{\partial f}{\partial x_i}(a)(x_i-a_i)+\frac{1}{2}\sum_{i,j=1}^n\frac{\partial^2 f}{\partial x_i\partial x_j}(a)(x_i-a_i)(x_j-a_j)+R_2(a,x)
    \end{equation*}
    where 
    \begin{equation*}
        \frac{R_2(a,x)}{\|x-a\|^2}\to 0 \text{ as } x\to a
    \end{equation*}
\end{defn}

















\begin{defn}[critical point]
    Let $f:U\subset\R^n\to\R$, a point $x_0\in U$ is a \textbf{critical point} of $f$ if either $f$ is not differentiable at $x_0$, or $Df(x_0)=0$. A critical point that is not a local extremum is called a saddle point.
\end{defn}

\begin{defn}[quadratic function]
    A function $g:\R^n\to\R$ is called a \textbf{quadratic function} if it is given by 
    \begin{equation*}
        g(h_1,\dots, h_n)=\sum_{i,j=1}^na_{ij}h_ih_j
    \end{equation*}
    where $(a_{ij})$ is an $n\times n$ matrix. We can also write $g$ as follows:
    \begin{equation*}
        g(h_1,\dots,h_n)=[h_1,\dots,h_n]\begin{bmatrix}
            a_{11}&\dots &a_{1n}\\
            \vdots&\ddots &\vdots\\
            a_{n_1}&\dots&a_{nn}
        \end{bmatrix}\begin{bmatrix}
            h_1\\
            \vdots\\
            h_n
        \end{bmatrix}
    \end{equation*}

\end{defn}

\begin{defn}[Hessian matrix]
    Let $f:U\subset\R^n\to\R$, and suppose all the second-order partial derivatives $\frac{\partial^2 f}{\partial x_i\partial x_j}$ exist, then the Hessian matrix of $f$ is the $n\times n$ matrix given by 
    \begin{equation*}
        Hf=\begin{bmatrix}
            \frac{\partial^2 f}{\partial x_1\partial x_1}&\dots&\frac{\partial^2f}{\partial x_1\partial x_n}\\
            \vdots&\ddots&\vdots\\
            \frac{\partial^2f}{\partial x_nx_1}&\dots&\frac{\partial^2f}{\partial x_n\partial x_n}
        \end{bmatrix}
    \end{equation*}
    The Hessian as a quadratic function is defined by 
    \begin{equation*}
        Hf(x)(h)=\frac{1}{2}\begin{bmatrix}
            h_1&\dots&h_n
        \end{bmatrix}Hf(x)\begin{bmatrix}
            h_1\\
            \dots\\
            h_n
        \end{bmatrix}
    \end{equation*}
    where $h=(h_1,\dots,h_n)$.
\end{defn}

\begin{defn}[degenerate/nondegenerate points]
    Let $f:U\subset\R^2\to\R$ be of $C^2$, let $(x_0,y_0)$ be a critical point. We define the \textbf{discriminant}, \boldmath $\mathcal{D}$,\unboldmath of the Hessian by
    \begin{equation*}
        \mathcal{D}=\det (Hf)=\left(\frac{\partial^2f}{\partial x^2}\right)\left(\frac{\partial^2f}{\partial y^2}\right)-\left(\frac{\partial^2f}{\partial x\partial y}\right)^2
    \end{equation*}
    If $\mathcal{D}\neq 0$, the critical point $(x_0,y_0)$ is called \textbf{nondegenerate}; if $\mathcal{D}=0$, the point $(x_0,y_0)$ is called \textbf{degenerate}.
\end{defn}

\begin{defn}[positive, negative-definite]
    A quadratic function $g:\R^n\to\R$ is called \textbf{positive-definite} if $g(h)\geq 0$ for all $h\in\R^n$ and $g(h)=0$ implies $h=0$. Similarly, $g$ is \textbf{negative-definite} if $g(h)\leq 0$ for all $h\in\R^n$ and $g(h)=0$ implies $h=0$. 
\end{defn}


\begin{defn}[global extremum]
    Let $f:A\to\R$ be a function defined on $A\subset\R^2$ or $A\subset\R^3$. A point $x_0\in A$ is said to be an \textbf{absolute maximum} if $f(x_0)\geq f(x)$ for all $x\in A$. Similarly, $x_0$ is an \textbf{absolute minimum} if $f(x_0)\leq f(x)$ for all $x\in A$. 
\end{defn}

\begin{defn}[bounded set]
    A set $A\subset\R^n$ is said to be \textbf{bounded} if there is a number $M>0$ such that $\|x\|\leq M$ for all $x\in A$. 
\end{defn}

\begin{defn}[acceleration]
    Let $c(t)$ be a path, the \textbf{acceleration} $a(t)$ of $c(t)$ is 
    \begin{equation*}
        a(t)=c''(t)
    \end{equation*}
\end{defn}



\begin{defn}[arc length]
    Let $c(t)=(x(t), y(t), z(t))$ be a path, then the length of the path in $\R^3$ from $t_0\leq t\leq t_1$ is 
    \begin{align*}
        L_{t_0\to t_1}(c)&=\int_{t_0}^{t_1}\left(x'(t)^2+y'(t)^2+z'(t)^2\right)^\frac{1}{2}dt\\
        &=\int_{t_0}^{t_1}\|c'(t)\|dt
    \end{align*}
    More generally, if $c(t)=(x_1(t), \dots, x_n(t))$ is a path in $\R^n$, then 
    \begin{equation*}
        L_{t_0\to t_1}(c)=\int_{t_0}^{t_1}\left(\sum_{i=1}^nx_i'(t)^2\right)^\frac{1}{2}dt
    \end{equation*}
\end{defn}


\begin{defn}[vector field]
    A vector field is a function $F:A\subset\R^n\to\R^n$ that assigns $x\in\R^n$ to another vector $F(x)\in\R^n$.
\end{defn}


\begin{defn}[flow line]
    If $F$ is a vector field, a \textbf{flow line} for $F$ is a path $c(t)$ such that 
    \begin{equation*}
        c'(t)=F(c(t))
    \end{equation*}
    Intuitively speaking, flow lines are the ``streamlines'' threading through vector fields.
\end{defn}



\begin{defn}[divergence]
    Let ${F}$ be a vector field in $\R^3$ $F=(F_1,F_2,F_3)$, the divergence of ${F}$ is the \textbf{scalar field} (assigns one number to an given point $(x, y, z)$), 
    \begin{equation*}
        \diverg F\coloneq\nabla\cdot F=\frac{\partial F_1}{\partial x}+\frac{\partial F_2}{\partial y}+\frac{\partial F_3}{\partial z}
    \end{equation*}
    More generally, if $F=(F_1,\dots, F_n)$ is a vector field on $\R^n$, its divergence is 
    \begin{equation*}
        \diverg F=\sum_{i=1}^n\frac{\partial F_i}{\partial x_i}=\frac{\partial F_1}{\partial x_1}+\dots+\frac{\partial F_n}{\partial x_n}
    \end{equation*}
\end{defn}


\begin{defn}[curl]
    Let $F$ be a vector field in $\R^3$, writing $F=(F_1,F_2,F_3)$, the \textbf{curl} of $F$ is the vector field 
    \begin{equation*}
        \curl F\coloneq\nabla\times F=\det\begin{bmatrix}
            i&j&k\\
            \frac{\partial}{\partial x}&\frac{\partial}{\partial y}&\frac{\partial}{\partial z}\\
            F_1&F_2&F_3
        \end{bmatrix}
    \end{equation*}
    If $\curl F=0$, then we say the vector field is \textbf{irrotational}.
\end{defn}



\begin{defn}
    We say a region $D\subset\R^2$ is \boldmath$y$\unboldmath-\textbf{simple} if there are continuous functions $\phi_1,\phi_2$ such that $D$ is the set of pionts $(x,y)$ satisfying
    \begin{equation*}
        x\in[a,b], \quad  \phi_1(x)\leq y\leq\phi_2(x)
    \end{equation*}
    Similarly, we define $D$ to be \boldmath$x$\textbf{-simple} \unboldmath if there are continuous $\psi_1, \psi_2$ such that $D$ is the set of points $(x,y)$ satisfying 
    \begin{equation*}
        y\in[c,d], \quad \psi_1(y)\leq x\leq\psi_2(y)
    \end{equation*}
    A \textbf{simple} region is one that is both $x$- and $y$-simple.
\end{defn}



\begin{defn}[injective, surjective]
    Let $T:\R^2\to\R^2$ be a map, we say $T$ is \textbf{injective}, or \textbf{one-to-one}, on $D^*$, if for $x,y\in D^*$
    \begin{equation*}
        Tx=Ty
    \end{equation*}
    implies
    \begin{equation*}
        x=y.
    \end{equation*}
    We say $T$ is \textbf{surjective}, or \textbf{onto} $D$, if for all $y\in D$, there exists $x$ in the domain of $T$ such that 
    \begin{equation*}
        Tx=y
    \end{equation*}
    If $T$ is both injective and surjective, then we say $T$ is \textbf{bijective}.
\end{defn}


\begin{defn}[Jacobian Determinant]
    Let $T:D^*\subset\R^2\to\R^2$ be of $C^1$ defined by 
    \begin{equation*}
        T:\begin{pmatrix}
            u\\
            v
        \end{pmatrix}\mapsto \begin{pmatrix}
            x(u,v)\\
            y(u,v)
        \end{pmatrix}
    \end{equation*}
    The \textbf{Jacobian determinant} of $T$, denoted as $\frac{\partial(x,y)}{\partial(u,v)}$ is the determinant of the matrix $DT(u,v)$:
    \begin{equation*}
        \frac{\partial(x,y)}{\partial(u,v)}=\det\left|\begin{matrix}
            \frac{\partial x}{\partial u}& \frac{\partial x}{\partial v}\\
            \frac{\partial y}{\partial u}&\frac{\partial y}{\partial v}
        \end{matrix}\right|
    \end{equation*}
\end{defn}


\chapter{Theorem Review}



\begin{prop}[extremums are critical points]
    Let $f:U\subset\R^n\to\R$ be differentiable, where $U$ is open. If $x_0$ is a local extremum, then $Df(x_0)=0$. 
\end{prop}


\begin{prop}[extremum]
    Let $f\colon U\subset\R^n\to\R$ be in $C^3$, and $x_0$ is a critical point of $f$. If the Hessian $Hf(x_0)$ is positive-definite, then $x_0$ is a local minimum of $f$; if $Hf(x_0)$ is negative-definite, then $x_0$ is a local maximum.
\end{prop}




\begin{prop}[local minimum]
    Let $f(x,y)$ be of $C^2$, and $U$ is open in $\R^2$. A point $(x_0,y_0)$ is a strict local \textbf{minimum} of $f$ if the following conditions hold:
    \begin{enumerate}
        \item \begin{equation*}
            \frac{\partial f}{\partial x}(x_0,y_0)=\frac{\partial f}{\partial y}(x_0,y_0)=0
        \end{equation*}
        \item \begin{equation*}
            \mathcal{D}(x_0,y_0)>0
        \end{equation*}
        where $\mathcal{D}$ is the \textbf{discriminant} of the Hessian, defined by 
        \begin{equation*}
            \mathcal{D}=\det (Hf)=\left(\frac{\partial^2 f}{\partial x^2}\right)\left(\frac{\partial^2 f}{\partial y^2}\right)-\left(\frac{\partial^2 f}{\partial x\partial y}\right)^2
        \end{equation*}
        where $Hf$ is the $2\times 2$ Hessian matrix.
        \item  \begin{equation*}
            \frac{\partial^2f}{\partial x^2}(x_0,y_0)>0
        \end{equation*}
    \end{enumerate}
\end{prop}


\begin{prop}[local maximum]
    Let $f(x,y)$ be of $C^2$, and $U$ is open in $\R^2$. A point $(x_0,y_0)$ is a strict local \textbf{maximum} of $f$ if the following conditions hold:
    \begin{enumerate}
        \item \begin{equation*}
            \frac{\partial f}{\partial x}(x_0,y_0)=\frac{\partial f}{\partial y}(x_0,y_0)=0
        \end{equation*}
        \item \begin{equation*}
            \mathcal{D}(x_0,y_0)>0
        \end{equation*}
        where $\mathcal{D}$ is the {discriminant} of the Hessian, defined above.
        \item  \begin{equation*}
            \frac{\partial^2f}{\partial x^2}(x_0,y_0)<0
        \end{equation*}
    \end{enumerate}
\end{prop}

\begin{prop}[saddle points]
    Let $f(x,y):U\subset\R^2\to \R$ be of  $C^2$, if $\frac{\partial f}{\partial x}(x_0,y_0)=\frac{\partial f}{\partial y}(x_0,y_0)=0$, and $\mathcal{D}(x_0,y_0)<0$, where $\mathcal{D}$ is the discriminant, then the critical point $(x_0,y_0)$ is a saddle point, i.e., neither a maximum or a minimum.
\end{prop}


\begin{prop}[continuous functions attain extremum on closed bounded sets]
    Let $f:D\to\R$ be continuous, where $D$ is closed and bounded in $\R^n$. Then $f$ assumes its absolute maximum and absolute minimum values at some point $x_0, x_1\in D$.
\end{prop}









\begin{prop}
    Let $f$ be constrained to a surface $S$, if $f$ has a max or a min at $x_0$, then $\nabla f(x_0)$ is perpendicular to $S$ at $x_0$.    
\end{prop}

\begin{prop}[Lagrange]
    Suppose that $f:U\subset\R^n\to\R$ and $g:U\subset\R^n\to\R$ are $C^1$ functions. Let $x_0\in U$ and $g(x_0)=c$, and let $S$ be the level set for $g$ at $c$, i.e., $S=\{x: g(x)=c\}$. Assume $\nabla g(x_0)\neq 0$, then if $f$ has a local maximum or minimum on $S$ at $x_0$, then there exists some real number $\lambda$ such that
    \begin{equation*}
        \nabla f(x_0)=\lambda \nabla g(x_0)
    \end{equation*}
\end{prop}


\begin{prop}[Bordered Hessian]
    Let $f: U\subset\R^2\to\R$ and $g:U\subset\R^2\to\R$ be smooth functions. Let $x_0\in U, g(x_0)=c$, and let $S$ be the level curve of $g$ with value $c$. Assume that $\nabla g(x_0)\neq 0$ and that there exists a real number $\lambda$ such that 
    \begin{equation*}
        \nabla f(x_0)=\lambda\nabla g(x_0)
    \end{equation*}
    Let $h=f-\lambda g$ and the bordered Hessian determinant is defined by 
    \begin{equation*}
        |\bar{H}|=\det\left|\begin{matrix}
            0&-\frac{\partial g}{\partial x}&-\frac{\partial g}{\partial y}\\
            -\frac{\partial g}{\partial x} &\frac{\partial^2h}{\partial x^2}&\frac{\partial^2h}{\partial x\partial y}\\
            -\frac{\partial g}{\partial y} &\frac{\partial^2h}{\partial x\partial y}&\frac{\partial^2h}{\partial y^2}
        \end{matrix} \right|
    \end{equation*}
    For $f$ restricted to the curve $S$, 
    \begin{enumerate}
        \item If  $|\bar{H}|>0$, then $x_0$ is a local max.
        \item If $|\bar{H}|<0$, then $x_0$ is a local min.
        \item If $|\bar{H}|=0$, then it is inconclusive.
    \end{enumerate}
\end{prop}


\begin{prop}[Newton's Second Law]
    Let $F$ be the force acting on a particle of mass $m$, then 
    \begin{equation*}
        F=ma
    \end{equation*}
    where $a$ is the acceleration.
\end{prop}



\begin{prop}[gradient is irrotational]
    Let $f\in C^2$, viewing $\nabla f$ as a vector field, then
    \begin{equation*}
        \nabla\times(\nabla f)=0
    \end{equation*}    
\end{prop}

\begin{prop}[divergence of a curl vanishes]
    For any $C^2$ vector field $F$, 
    \begin{equation*}
        \nabla\cdot(\nabla\times F)=0
    \end{equation*}
\end{prop}



\begin{prop}[Fubini's Theorem for rectangles]
    Let $f$ be a continuous function on a rectangular domain $R=[a,b]\times[c,d]$, then 
    \begin{equation*}
        \int_a^b\int_c^df(x,y)dydx=\int_c^d\int_a^bf(x,y)dxdy
    \end{equation*}
\end{prop}

\begin{prop}[Fubini's Theorem for general regions]
    Suppose $D$ is a set of points $(x,y)$ such that $y\in [c,d]$ and $\psi_1(y)\leq x\leq\psi_2(y)$, and similarly for $x\in [a,b]$, $\varphi_1(x)\leq y\leq\varphi_2(x)$. If $f$ is continuous on $D$, then 
    \begin{equation*}
        \iint_Df(x,y)dA=\int_a^b\int_{\varphi_1(x)}^{\varphi_2(x)}f(x,y)dydx=\int_c^d\int_{\psi_1(y)}^{\psi_2(y)}f(x,y)dxdy
    \end{equation*}
\end{prop}


\begin{prop}
    We have the following identities regarding divergence and curl: 
    \begin{enumerate}
        \item $\nabla(f+g)=\nabla f+\nabla g$.
        \item $\nabla (cf)=c\nabla f$, for constant $c$.
        \item $\nabla (fg)=f\nabla g+g\nabla f$.
        \item $\nabla (f/g)=(g\nabla f-f\nabla g)/g^2$, at points $x$ where $g(x)\neq 0$.
        \item $\diverg (F+G)=\diverg F+\diverg G$.
        \item $\curl(F+G)=\curl F+\curl G$.
        \item $\diverg(fF)=f\diverg F+F\cdot\nabla f$.
        \item $\diverg (F\times G)=G\cdot\curl F-F\cdot\curl G$.
        \item $\diverg\curl F=0$.
        \item $\curl(fF)=f\curl F+\nabla f\times F$.
        \item $\curl\nabla f=0$.
        \item $\nabla^2(fg)=f\nabla^2g+g\nabla^2f+2(\nabla f\cdot \nabla g)$.
        \item $\diverg(\nabla f\times\nabla g)=0$.
        \item $\diverg(f\nabla g-g\nabla f)=f\nabla^2g-g\nabla^2f$.
    \end{enumerate}
\end{prop}



\begin{prop}[integrability]
    For different assumptions on $f$, we have the following integrability results:
    \begin{enumerate}
        \item Let $f$ be continuous and defined on a closed rectangle $R$, then $f$ is integrable over $R$.
        \item Let $f:R\to\R$ be a bounded function on $R$ and suppose the set of points where $f$ is discontinuous lies on a finite union of graphs of continuous functions, then $f$ is integrable over $R$.
    \end{enumerate}
\end{prop}


\begin{prop}[Fubini's Theorem for rectangles]
    For different assumptions on $f$, we have the following Fubini's theorem results:
    \begin{enumerate}
        \item   Let $f$ be a continuous function on a rectangular domain $R=[a,b]\times[c,d]$, then 
        \begin{equation*}
            \int_a^b\int_c^df(x,y)dydx=\int_c^d\int_a^bf(x,y)dxdy=\iint_Rf(x,y)dA
        \end{equation*}
        \item Let $f$ be bounded with domain $R=[a,b]\times[c,d]$ and the discontinuities of $f$ lie on a finite union of graphs of continuous functions. If the integral $\int_a^bfdy$ exists for each $x\in[a,b]$, then 
        \begin{equation*}
            \int_a^b\left(\int_c^df(x,y)dy\right)dx
        \end{equation*}
        exists and 
        \begin{equation*}
            \int_a^b\int_c^df(x,y)dydx=\iint_Rf(x,y)dA
        \end{equation*}
        Similar results hold if $\int_a^bfdx$ exists for each $y\in[c,d]$. If both hold simultaneously, then 
        \begin{equation*}
            \int_a^b\left(\int_c^df(x,y)dy\right)dx=\int_c^d\left(\int_a^bf(x,y)dx\right)dy=\iint_Rf(x,y)dA
        \end{equation*}
    \end{enumerate}

  
\end{prop}



\begin{prop}[Fubini's Theorem for general regions]
    Suppose $D$ is a set of points $(x,y)$ such that $y\in [c,d]$ and $\psi_1(y)\leq x\leq\psi_2(y)$, and similarly for $x\in [a,b]$, $\varphi_1(x)\leq y\leq\varphi_2(x)$. If $f$ is continuous on $D$, then 
    \begin{equation*}
        \iint_Df(x,y)dA=\int_a^b\int_{\varphi_1(x)}^{\varphi_2(x)}f(x,y)dydx=\int_c^d\int_{\psi_1(y)}^{\psi_2(y)}f(x,y)dxdy
    \end{equation*}
\end{prop}



\begin{prop}[simple-regions]
    If $D$ is a $x$-simple region with $y\in[c,d], \psi_1(y)\leq x\leq\psi_2(y)$, and if $f$ is continuous on $D$, then 
    \begin{equation*}
        \iint_Df(x,y)dA=\int_c^d\left(\int_{\psi_1(y)}^{\psi_2(y)}f(x,y)dx\right)dy
    \end{equation*}
    Similarly, if $D$ is $y$-simple, then
    \begin{equation*}
        \iint_Df(x,y)dA=\int_a^b\int_{\phi_1(x)}^{\phi_2(x)}f(x,y)dydx
    \end{equation*}
    If $D$ is simple, then the two expressions above are equal.
\end{prop}
For example, the area of a $x$-simple region $D$ can be computed as 
    \begin{equation*}
        \iint_DdA=\int_c^d\psi_2(y)-\psi_1(y)dy
    \end{equation*}

\begin{prop}
    Let $A$ be $2\times 2$ matrix with $\det(A)\neq 0$ and let $T:\R^2\to\R^2$ be the linear map $Tx=Ax$. Then $T$ transforms parallelograms into parallelograms and vertices into vertices.
\end{prop}

\begin{prop}
    Let $T:\R^n\to\R^n$ be a linear map, i.e., there exists $n\times n$ matrix $A$ such that $Tx=Ax$, then $T$ is injective iff surjective iff $\det(A)\neq 0$.
\end{prop}


\begin{thm}[change of variables formula]
    Let $D,D^*$ be elementary regions in $\R^2$, suppose $T:D^*\to D$ is both one-to-one and onto. Then for any integral function $f:D\to\R$, the \textbf{change of variable formula} states
    \begin{equation*}
        \iint_Df(x,y)dxdy=\iint_{D^*}f(x(u,v),y(u,v))\left|\det(J)\right|dudv
    \end{equation*}
    where 
    \begin{equation*}
        \det(J)=\left|\frac{\partial(x,y)}{\partial(u,v)}\right|
    \end{equation*}
    is the Jacobian determinant. 
\end{thm}

\begin{prop}[change of variables-polar coordinates]
    As a corollary to the theorem above, we have the following change of variables formula for polar coordinates:
    \begin{equation*}
        \iint_Df(x,y)dxdy=\iint_{D^*}f(r\cos\theta, r\sin\theta)rdrd\theta
    \end{equation*}
\end{prop}



\begin{prop}[change of variables-triple]
    Let $W,W^*$ be elementary regions in $\R^3$, and suppose $T:W^*\to W$ is bijective. Then the change of varialbes formula for triple integrals states:
    \begin{equation*}
        \iiint_Wf(x,y,z)dxdydz=\iiint_{W^*}f(x(u,v,w), y(u,v,w), z(u,v,w))\mid\det(J)\mid dudvdw
    \end{equation*}
    where 
    \begin{equation*}
        \det(J)=\det\left|\begin{matrix}
            \frac{\partial x}{\partial u}&\frac{\partial x}{\partial v}&\frac{\partial x}{\partial w}\\
            \frac{\partial y}{\partial u}&\frac{\partial y}{\partial v}&\frac{\partial y}{\partial w}\\
            \frac{\partial z}{\partial u}&\frac{\partial z}{\partial v}&\frac{\partial z}{\partial w}
        \end{matrix}\right|
    \end{equation*}
    is the Jacobian determinant.
\end{prop}


\begin{prop}[change of variables-triple cylindrical]
    As a corollary to the above, we have the following change of variables formula for cylindrical coordinates:
    \begin{equation*}
        \iiint_Wf(x,y,z)dxdydz=\iiint_{W^*}f(r\cos\theta, r\sin\theta, z)rdrd\theta dz
    \end{equation*}
    Recall cylinder cooridnates is setting up the following 
    \begin{equation*}
        x=r\cos\theta, y=r\sin\theta, z=z
    \end{equation*}
\end{prop}

\begin{prop}[change of variables-triple spherical]
    As a corollary to the above, we have the following change of variables formula for spherical coordinates:
    \begin{equation*}
        \iiint_Wf(x,y,z)dxdydz=\iiint_{W^*}f(\rho\sin\phi\cos\theta, \rho\sin\phi\sin\theta, \rho\cos\phi)\rho^2\sin\phi d\rho d\theta d\phi
    \end{equation*}
    Recall the spherical cooridnates is setting up the following 
    \begin{equation*}
        x=\rho\sin\phi\cos\theta, y=\rho\sin\phi\sin\theta, z=\rho\cos\phi
    \end{equation*}


\end{prop}



\chapter{Practice Problems}





























\end{document}